
\Slide{1}{关系的起点}
二元关系比较两个对象。函数强调输入输出，关系强调哪些有序对满足条件；一个输入可与多个对象有关。

\Slide{2}{知识地图}
集合、函数之后引入关系，再用关系描述大小、偏好和确界。实数的顺序公理会用到这一讲的语言。

\Slide{3}{关系与例子}
$R\subseteq S\times S$，$xRy$ 是 $(x,y)\in R$ 的简写。包含关系是子集之间的关系，而不是元素之间的属于关系；两者符号不可交换。

\Slide{4}{预序与偏序}
预序满足自反和传递；偏序额外要求反对称。反对称是“双向关系推出相等”，不是“永远不能双向”。例如通常的 $\le$ 在 $x=y$ 时允许双向。

\Slide{5}{严格与无差异}
$x\sim y$ 表示双向弱关系；$x\prec y$ 表示一个方向成立且反方向不成立。无差异未必是对象相同，例如两种不同消费带来同样评价。

\Slide{6}{完全性}
完全性要求任意两点可比较。坐标序中 $(1,0)$ 与 $(0,1)$ 不可比较，所以不是完全序。注意“不能证明谁更好”不是数学上的无差异。

\Slide{7}{关系分类图}
沿箭头增加条件会缩小适用对象。检查一种关系属于哪类，依次检查自反、传递、反对称、完全；否定一个性质只需找到违反它的具体对象。

\Slide{8}{极大与最大}
极大是不被严格压过，最大是压过所有其他对象。整除序里 $3,4,5$ 都极大，但没有最大。完全预序中极大即最大：任一点与它可比较，不能严格更大，就只能不超过它。

\Slide{9}{上界与确界}
$u$ 是上界要求每个 $x\in A$ 满足 $x\preceq u$；确界还要求 $u$ 不大于任何其他上界。上界通常不在集合中。例如 $(0,1)$ 的确界是1，没有最大值。

\Slide{10}{完备性预告}
非空有上界集合有确界是实数的关键性质，不是所有有序集合都满足。正有理数中平方小于2的数有上界，但在有理数内没有确界。

\Slide{11}{区间}
区间包含任意两点之间的全部可比较点。实数中 $[a,b]$ 包含端点，$(a,b)$ 不含端点。写邻域时括号变化会影响端点是否属于集合。

\Slide{12}{乘积序}
$x\le y$ 表示逐坐标不大于；$x<y$ 表示逐坐标不大于且至少一处严格；$x\ll y$ 则要求每处严格。例 $(0,1)<(1,1)$ 但不是 $\ll$。

\Slide{13}{单调映射}
序保持是 $x\preceq_S y\Rightarrow f(x)\preceq_T f(y)$。证明须从任意可比较的输入开始，不能只测几个值。严格递增、递减与弱递增的条件应分别写明。

\Slide{14}{有序对的正式定义}
关系是集合，因此能取补、逆、交等。反关系把 $(x,y)$ 换成 $(y,x)$，补关系表示原判断不成立。以后“负传递”就是补关系的传递性。

\Slide{15}{七种性质逐项读}
自反与非自反不互为一般否定：否定自反只要求存在一个点不与自身有关。非对称推出反对称，但反对称允许对角点。做题把每个定义的全称量词写出来。

\Slide{16}{严格偏好的要求}
本课程用非对称与负传递刻画严格偏好。负传递是 $x\not\succ y,\ y\not\succ z\Rightarrow x\not\succ z$，不是普通传递换一个名字。

\Slide{17}{负传递的等价形式}
原式的逆否命题是 $x\succ z\Rightarrow(x\succ y\text{ 或 }y\succ z)$，对每个中间点 $y$ 成立。推导用到 $\neg(P\land Q)=\neg P\lor\neg Q$，注意“或”允许两者同时成立。

\Slide{18}{严格偏好定义}
这是后面所有偏好推论的假设；不把它当作现实中自动满足的规律。出现循环选择就可能违反这些条件，数学模型是否适合现实需另行判断。

\Slide{19}{传递与无循环的证明}
非对称排除 $x\succ x$。若 $x\succ y\succ z$，由上页等价式，$x\succ y$ 推出 $x\succ z$ 或 $z\succ y$；后者与 $y\succ z$ 矛盾，故传递。沿有限循环反复传递得到 $x\succ x$，矛盾。

\Slide{20}{从严格定义弱偏好}
定义 $x\succeq y\iff\neg(y\succ x)$，无差异表示两个严格方向都不成立。方向反转是易错点：弱偏好是“不严格差于”，不是“不严格好于”。

\Slide{21}{三个互斥情形}
非对称排除同时双向严格，余下恰为无差异。弱偏好传递由负传递得到：$\neg(y\succ x)$ 与 $\neg(z\succ y)$ 推出 $\neg(z\succ x)$。无差异的传递来自两条弱偏好链分别传递。

\Slide{22}{命题证明展开}
若 $x\succ y$ 且 $y\sim z$，由负传递等价式 $x\succ z$ 或 $z\succ y$，后一项被无差异排除。若还 $z\succ w$，再用传递得 $x\succ w$。每个被排除的分支都应写出所违反的假设。

\Slide{23}{两种原始描述}
本页解释以严格偏好为基础的建模选择。弱偏好和无差异是用定义导出的，不应再任意赋值，否则可能互相矛盾。

\Slide{24}{弱偏好定义}
完全性包括与自身比较，因而推出自反。传递要求任何三点的两段比较能够连接。没有完全性时，$\neg(y\succeq x)$ 不足以推出 $x\succeq y$。

\Slide{25}{反向构造}
给定完全传递的 $\succeq$，定义严格关系为弱关系成立且反向不成立；完全性使其等价于单写反向不成立。无差异仍为双向弱关系。

\Slide{26}{构造的性质}
严格关系非对称：若 $x\succ y$，则 $x\succeq y$，故不能同时 $y\succ x$。负传递：$x\not\succ y$ 意味 $y\succeq x$，再与 $z\succeq y$ 传递即得 $x\not\succ z$。无差异为等价关系，由自反、交换方向和双链传递得到。

\Slide{27}{原页无差异公式纠错}
原图第二项应为 $x_2\succeq' x_1$，不能是 $x_2\succeq' x_2$。后者总真，会把无差异误写成单向弱偏好。修正后上下两种出发方式互相还原。

\Slide{28}{还原弱关系}
原构造给出 $x\succeq y\iff\neg(y\succ' x)\iff\neg\neg(x\succeq' y)\iff x\succeq' y$。严格关系和无差异也由相同定义还原。这里的双重否定是逻辑步骤，不是经验判断。

\Slide{29}{本讲复习}
分别从严格和弱偏好出发重写定义，检查最终回到原关系。自测：用效用 $u$ 定义 $x\succeq y\iff u(x)\ge u(y)$；实数顺序的完全和传递立即给出完全预序，但 $u(x)=u(y)$ 不推出 $x=y$，所以一般不是偏序。
